% ===========================================================================
%  Chapitre 3 — Étude et représentation graphique des fonctions
%  PE 2026, composante ANALYSE
% ===========================================================================
\bmtchapitre{Étude et représentation graphique des fonctions}
  {Étudier le comportement asymptotique d'une courbe, repérer points angulés et
   points d'inflexion, situer une courbe par rapport à ses asymptotes, et
   résoudre graphiquement équations et inéquations.}
  {Analyse \textbullet\ environ 8 heures}

Vous savez maintenant calculer une limite et une dérivée. Ce chapitre ne
présente presque aucune notion nouvelle : il apprend à \emph{les faire tenir
ensemble} sur une feuille, jusqu'à produire une courbe juste.

C'est un chapitre de méthode, et c'est aussi celui qui rapporte le plus de
points au baccalauréat. Le tracé final vaut souvent un point et demi à lui
seul, et il ne s'improvise pas : chaque trait doit avoir été justifié avant
d'être posé. Une courbe qui ignore une asymptote qu'on vient de démontrer,
c'est un point perdu deux fois — une fois pour le dessin, une fois pour la
cohérence.

% ---------------------------------------------------------------------------
\begin{revision}
Rien de neuf ici : on vérifie que les deux premiers chapitres sont en place.

\begin{enumerate}[label=\textbf{\arabic*.}, leftmargin=*, itemsep=0.35em]
  \item Déterminer $\limpinf \dfrac{3x^{2}-x+1}{x^{2}+2}$.
  \item Que signifie « $f$ est continue en $a$ » ?
  \item Comment reconnaît-on un point angulé sur une courbe ?
  \item Quel lien y a-t-il entre le signe de $f''$ et la forme de la courbe ?
  \item Donner l'ensemble de définition de $x \mapsto \dfrac{x^{2}+1}{x-3}$.
  \item Effectuer la division euclidienne de $x^{2}+1$ par $x-3$.
\end{enumerate}
\end{revision}

\begin{automatismes}
Sans calculatrice, sans brouillon.

\begin{enumerate}[label=\textbf{\alph*)}, leftmargin=*, itemsep=0.25em]
  \item $\limpinf \dfrac{2x+1}{x-4}$
  \item $\displaystyle\lim_{x \to 4^{+}} \dfrac{2x+1}{x-4}$
  \item $\limpinf \dfrac{x^{2}}{x+1}$
  \item Dérivée de $x \mapsto \dfrac{x^{2}}{x+1}$
  \item Signe de $x^{2}-4$
  \item $\limminf (x^{3}-x)$
  \item Résoudre $\dfrac{1}{x-2} > 0$
  \item Le quotient $\dfrac{x^{2}}{x+1}$ vaut-il $x-1+\dfrac{1}{x+1}$ ?
\end{enumerate}
\end{automatismes}

\begin{corrige}[automatismes]
a) $2$ \quad b) $+\infty$ \quad c) $+\infty$ \quad
d) $\dfrac{x^{2}+2x}{(x+1)^{2}}$ \quad e) négatif sur $\intervalleo{-2}{2}$,
positif ailleurs \quad f) $-\infty$ \quad g) $x > 2$ \quad
h) oui, c'est la division euclidienne de $x^{2}$ par $x+1$.
\end{corrige}

% ===========================================================================
\section{Les branches infinies}

Lorsque $x$ ou $f(x)$ s'échappe vers l'infini, la courbe part au loin. Reste à
savoir \emph{comment} : en suivant une droite, ou en s'en écartant. C'est toute
la question des branches infinies.

\subsection{Asymptotes verticales et horizontales}

\begin{definition}[asymptote verticale]
Si $\displaystyle\lim_{x \to a} f(x) = \pm\infty$ — la limite pouvant n'être
qu'à gauche ou qu'à droite — alors la droite d'équation $x = a$ est
\textbf{asymptote verticale} à la courbe de $f$.
\end{definition}

\begin{definition}[asymptote horizontale]
Si $\displaystyle\lim_{x \to +\infty} f(x) = \ell$ avec $\ell$ réel, alors la
droite d'équation $y = \ell$ est \textbf{asymptote horizontale} à la courbe en
$+\infty$. Même définition en $-\infty$.
\end{definition}

\subsection{Asymptotes obliques}

\begin{definition}[asymptote oblique]
La droite d'équation $y = ax+b$, avec $a \neq 0$, est \textbf{asymptote
oblique} à la courbe de $f$ en $+\infty$ lorsque
\[
  \limpinf \bigl[f(x) - (ax+b)\bigr] = 0 .
\]
\end{definition}

\begin{methode}{trouver une asymptote oblique}
Deux chemins, selon la forme de $f$.

\smallskip
\emph{Si $f$ est une fraction rationnelle}, le plus rapide est la division
euclidienne du numérateur par le dénominateur. Écrire
$f(x) = ax+b+\dfrac{r(x)}{q(x)}$ avec $\deg r < \deg q$ : la partie entière
$ax+b$ est l'asymptote, et le reste tend vers $0$.

\smallskip
\emph{Sinon}, on cherche $a$ puis $b$ séparément :
\[
  a = \limpinf \frac{f(x)}{x}, \qquad b = \limpinf \bigl[f(x)-ax\bigr].
\]
Si ces deux limites existent et sont finies avec $a \neq 0$, la droite
$y = ax+b$ est asymptote.
\end{methode}

\begin{propriete}[position relative de la courbe et de l'asymptote]
Le signe de la différence $d(x) = f(x)-(ax+b)$ donne la position :
\begin{itemize}[leftmargin=1.4em, itemsep=0.15em]
  \item $d(x) > 0$ : la courbe est \emph{au-dessus} de l'asymptote ;
  \item $d(x) < 0$ : la courbe est \emph{au-dessous} ;
  \item $d(x) = 0$ : la courbe \emph{coupe} l'asymptote.
\end{itemize}
\end{propriete}

\begin{visualisation}[la courbe rejoint sa droite]
\begin{center}
\begin{tikzpicture}
\begin{axis}[
  width = 11cm, height = 6.2cm,
  axis lines = middle, xlabel = {$x$},
  xmin = -5.6, xmax = 5.6, ymin = -6.5, ymax = 6.5,
  xtick = \empty, ytick = \empty,
  axis line style = {bmtGris}, clip = true,
]
  % f(x) = x + 1/x  : asymptote y = x, asymptote verticale x = 0
  \addplot[bmtIndigo, line width=1.3pt, domain=0.18:5.4, samples=200] {x + 1/x};
  \addplot[bmtIndigo, line width=1.3pt, domain=-5.4:-0.18, samples=200] {x + 1/x};
  \addplot[bmtLaterite, dashed, line width=0.9pt, domain=-5.4:5.4, samples=2] {x};
  \node[bmtLaterite, font=\footnotesize, anchor=west] at (axis cs:4.0,3.1) {$y=x$};
  \node[bmtGris, font=\footnotesize, anchor=west] at (axis cs:0.25,-5.4) {$x=0$};
\end{axis}
\end{tikzpicture}
\end{center}

Regardez la branche de droite : plus $x$ grandit, plus la courbe se colle à la
droite rouge sans jamais la toucher, en restant au-dessus. La branche de
gauche fait la même chose, mais par en dessous. C'est exactement ce que dit le
signe de la différence : ici $f(x)-x = \dfrac{1}{x}$, positif à droite,
négatif à gauche. La droite verticale, elle, marque l'endroit où la fonction
cesse d'être définie.
\end{visualisation}

\subsection{Quand il n'y a pas d'asymptote}

\begin{definition}[direction asymptotique]
Supposons $\limpinf f(x) = \pm\infty$ et posons $a = \limpinf \dfrac{f(x)}{x}$.
\begin{itemize}[leftmargin=1.4em, itemsep=0.2em]
  \item Si $a = \pm\infty$, la courbe admet une \textbf{branche parabolique de
        direction $(Oy)$}.
  \item Si $a = 0$, elle admet une \textbf{branche parabolique de direction
        $(Ox)$}.
  \item Si $a$ est un réel non nul mais que $f(x)-ax$ tend vers l'infini, la
        courbe admet une \textbf{direction asymptotique} de coefficient $a$,
        sans asymptote.
\end{itemize}
\end{definition}

\begin{exemple}[trois comportements en $+\infty$]
\emph{La fonction $x \mapsto x^{2}$.} Ici $\dfrac{f(x)}{x} = x \to +\infty$ :
branche parabolique de direction $(Oy)$. La courbe monte plus vite que
n'importe quelle droite.

\smallskip
\emph{La fonction $x \mapsto \sqrt{x}$.} Ici
$\dfrac{f(x)}{x} = \dfrac{1}{\sqrt{x}} \to 0$ : branche parabolique de
direction $(Ox)$. La courbe monte, mais de moins en moins vite.

\smallskip
\emph{La fonction $x \mapsto x + \sqrt{x}$.} Ici $\dfrac{f(x)}{x} \to 1$, mais
$f(x)-x = \sqrt{x} \to +\infty$ : direction asymptotique de coefficient $1$,
et pourtant aucune asymptote. La courbe suit la direction de la première
bissectrice tout en s'en éloignant indéfiniment.
\end{exemple}

\begin{erreurclassique}
Trouver $a = \limpinf \dfrac{f(x)}{x}$ fini ne suffit pas à conclure à une
asymptote. Il faut ensuite vérifier que $f(x)-ax$ admet une limite
\emph{finie}. Le troisième cas de l'exemple ci-dessus est précisément le
piège : beaucoup de copies s'arrêtent au calcul de $a$ et annoncent une
asymptote qui n'existe pas.
\end{erreurclassique}

\begin{entrainement}[branches infinies]
\exo[\niveau{1}] Déterminer les asymptotes de la courbe de
$f(x) = \dfrac{3x-1}{x+2}$.

\exo[\niveau{2}] Soit $f(x) = \dfrac{x^{2}-x+1}{x-1}$. Montrer par division
euclidienne que la droite $y = x$ est asymptote à la courbe, et étudier la
position relative.

\exo[\niveau{2}] Étudier les branches infinies de $f(x) = \dfrac{x^{3}}{x^{2}-1}$.

\exo[\niveau{3}] Soit $f(x) = \sqrt{x^{2}+x}$ définie pour $x \geq 0$. Montrer
que la droite $y = x+\dfrac{1}{2}$ est asymptote à la courbe en $+\infty$.
\end{entrainement}

\begin{corrige}
\textbf{1.} $x = -2$ est asymptote verticale ; $\lim_{x \to \pm\infty} f(x) = 3$ donc $y = 3$ est
asymptote horizontale en $\pm\infty$.

\textbf{2.} $\dfrac{x^{2}-x+1}{x-1} = x + \dfrac{1}{x-1}$. Le reste tend vers
$0$ : $y = x$ est asymptote. Son signe est celui de $x-1$ : la courbe est
au-dessus pour $x>1$, au-dessous pour $x<1$.

\textbf{3.} $\dfrac{x^{3}}{x^{2}-1} = x + \dfrac{x}{x^{2}-1}$, donc $y=x$ est
asymptote oblique. De plus $x = 1$ et $x = -1$ sont asymptotes verticales.

\textbf{4.} $f(x)-x = \dfrac{(x^{2}+x)-x^{2}}{\sqrt{x^{2}+x}+x}
= \dfrac{x}{\sqrt{x^{2}+x}+x}$, qui tend vers $\dfrac{1}{2}$ en $+\infty$
(diviser numérateur et dénominateur par $x$). Donc $f(x)-\left(x+\tfrac12\right) \to 0$.
\end{corrige}

% ===========================================================================
\section{Ce que la courbe montre en un point}

Trois accidents peuvent se produire en un point, et le vocabulaire du
programme les distingue soigneusement.

\begin{definition}[trois points remarquables]
Soit $a$ un réel et $(\Cf)$ la courbe de $f$.
\begin{itemize}[leftmargin=1.4em, itemsep=0.25em]
  \item \textbf{Point de discontinuité} : $f$ n'est pas continue en $a$. La
        courbe est interrompue ; il faut lever le crayon.
  \item \textbf{Point angulé} : $f$ est continue en $a$, dérivable à gauche et
        à droite, mais $f'_{g}(a) \neq f'_{d}(a)$. Deux demi-tangentes,
        aucune tangente.
  \item \textbf{Point d'inflexion} : $f''$ s'annule en $a$ \emph{en changeant
        de signe}. La tangente traverse la courbe, qui passe de concave à
        convexe ou l'inverse.
\end{itemize}
\end{definition}

\begin{remarque}
Un quatrième cas mérite d'être connu : lorsque le taux d'accroissement tend
vers $\pm\infty$, la courbe admet une \emph{tangente verticale}. La fonction
n'y est pas dérivable, mais la courbe n'y présente aucun angle — elle se
redresse simplement à la verticale, comme $x \mapsto \sqrt{x}$ en $0$.
\end{remarque}

% ===========================================================================
\section{Le plan d'étude, une fois pour toutes}

\begin{methode}{étudier complètement une fonction}
Le même enchaînement sert pour toutes les fonctions de l'année. Suivez-le dans
l'ordre : chaque étape utilise la précédente.

\begin{enumerate}[label=\textbf{\arabic*.}, leftmargin=*, itemsep=0.28em]
  \item \textbf{Ensemble de définition.} Écarter les valeurs interdites
        (dénominateur nul, radicande négatif).
  \item \textbf{Parité, périodicité.} Si $f$ est paire ou impaire, on réduit
        l'étude à une moitié de l'ensemble et l'on complète par symétrie.
  \item \textbf{Limites aux bornes.} Aux bornes ouvertes de l'ensemble de
        définition, sans en oublier une seule.
  \item \textbf{Branches infinies.} Asymptotes verticales, horizontales,
        obliques ; directions asymptotiques ; position relative.
  \item \textbf{Dérivée et signe.} Calculer $f'$, la factoriser, en étudier le
        signe.
  \item \textbf{Tableau de variation.} Y reporter limites, valeurs
        remarquables et flèches.
  \item \textbf{Points particuliers.} Intersections avec les axes, tangentes
        remarquables, points d'inflexion si la question le demande.
  \item \textbf{Tracé.} Placer d'abord les asymptotes, puis les points connus,
        puis la courbe — dans cet ordre, jamais l'inverse.
\end{enumerate}
\end{methode}

\begin{exemple}[une étude menée jusqu'au bout]
Étudions $f(x) = \dfrac{x^{2}-2x+2}{x-1}$.

\emph{Définition.} $\Df = \R \setminus \{1\}$.

\emph{Forme réduite.} La division euclidienne donne
$f(x) = x-1+\dfrac{1}{x-1}$.

\emph{Limites.} En $\pm\infty$, $f(x) \to \pm\infty$. En $1^{-}$ :
$\dfrac{1}{x-1} \to -\infty$ donc $f(x) \to -\infty$ ; en $1^{+}$,
$f(x) \to +\infty$.

\emph{Branches infinies.} $x=1$ est asymptote verticale. Comme
$f(x)-(x-1) = \dfrac{1}{x-1} \to 0$, la droite $y = x-1$ est asymptote
oblique aux deux infinis. Son signe est celui de $x-1$ : la courbe est
au-dessus pour $x>1$, au-dessous pour $x<1$.

\emph{Dérivée.} $f'(x) = 1 - \dfrac{1}{(x-1)^{2}} = \dfrac{(x-1)^{2}-1}{(x-1)^{2}}
= \dfrac{x(x-2)}{(x-1)^{2}}$.

Le signe de $f'$ est celui de $x(x-2)$ : positif à l'extérieur de
$\intervalle{0}{2}$, négatif à l'intérieur.

\emph{Variations.} $f$ croît sur $\left]-\infty ; 0\right]$, décroît sur
$\intervallefo{0}{1}$ puis sur $\intervalleof{1}{2}$, croît sur
$\left[2 ; +\infty\right[$. Les extremums locaux valent $f(0) = -2$ et
$f(2) = 2$.

\emph{Tracé.} On place les deux asymptotes, les points $(0;-2)$ et $(2;2)$,
puis on relie en respectant les variations et les positions établies.
\end{exemple}

\begin{visualisation}[l'étude précédente, achevée]
\begin{center}
\begin{tikzpicture}
\begin{axis}[
  width = 11cm, height = 6.4cm,
  axis lines = middle, xlabel = {$x$},
  xmin = -3.6, xmax = 5.6, ymin = -6.5, ymax = 6.5,
  xtick = {0,2}, ytick = {-2,2},
  axis line style = {bmtGris},
  tick label style = {font=\footnotesize, color=bmtGris}, clip = true,
]
  \addplot[bmtIndigo, line width=1.3pt, domain=1.12:5.4, samples=220] {x - 1 + 1/(x-1)};
  \addplot[bmtIndigo, line width=1.3pt, domain=-3.4:0.88, samples=220] {x - 1 + 1/(x-1)};
  \addplot[bmtLaterite, dashed, line width=0.9pt, domain=-3.4:5.4, samples=2] {x-1};
  \addplot[bmtVetiver, dashed, line width=0.9pt] coordinates {(1,-6.5) (1,6.5)};
  \filldraw[bmtIndigo] (axis cs:0,-2) circle (1.8pt);
  \filldraw[bmtIndigo] (axis cs:2,2) circle (1.8pt);
  \node[bmtLaterite, font=\footnotesize, anchor=west] at (axis cs:4.2,2.6) {$y=x-1$};
  \node[bmtVetiver, font=\footnotesize, anchor=south west] at (axis cs:1.1,-6.2) {$x=1$};
\end{axis}
\end{tikzpicture}
\end{center}

Comparez ce dessin à l'étude : les deux droites en pointillé ont été tracées
\emph{avant} la courbe, et c'est ce qui rend le tracé possible. Le maximum
local en $(0;-2)$ est sous l'axe des abscisses, le minimum local en $(2;2)$
au-dessus : la courbe ne coupe donc jamais l'axe des abscisses, ce que
confirme le discriminant négatif de $x^{2}-2x+2$.
\end{visualisation}

\begin{entrainement}[études complètes]
\exo[\niveau{2}] Étudier et représenter $f(x) = \dfrac{2x^{2}-3x}{x-1}$.

\exo[\niveau{2}] Étudier et représenter $f(x) = \dfrac{x^{2}+1}{x}$, après
avoir montré que $f$ est impaire.

\exo[\niveau{3}] Soit $f(x) = \dfrac{x^{2}+ax+b}{x-1}$. Déterminer $a$ et $b$
pour que la courbe admette la droite $y = x+2$ comme asymptote oblique et
passe par le point de coordonnées $(0;3)$.
\end{entrainement}

\begin{corrige}
\textbf{1.} $f(x) = 2x-1-\dfrac{1}{x-1}$ ; asymptotes $x=1$ et $y=2x-1$ ;
$f'(x) = \dfrac{2x^{2}-4x+3}{(x-1)^{2}}$, de discriminant négatif donc $f'>0$ :
$f$ est croissante sur chacun des deux intervalles.

\textbf{2.} $f(-x) = -f(x)$ : impaire, courbe symétrique par rapport à
l'origine. $f(x) = x+\dfrac{1}{x}$ ; asymptotes $x=0$ et $y=x$ ;
$f'(x) = 1-\dfrac{1}{x^{2}}$, nulle en $\pm1$ : minimum $2$ en $1$, maximum
$-2$ en $-1$.

\textbf{3.} La division donne
$f(x) = x+(a+1)+\dfrac{a+b+1}{x-1}$. L'asymptote $y = x+2$ impose $a+1 = 2$,
soit $a = 1$. Le passage par $(0;3)$ donne $\dfrac{b}{-1} = 3$, soit
$b = -3$. On vérifie que $a+b+1 = -1 \neq 0$, donc l'asymptote est bien
oblique.
\end{corrige}

% ===========================================================================
\section{Lire et résoudre sur un graphique}

Une courbe donnée remplace parfois tout un calcul. Le programme demande de
savoir résoudre graphiquement quatre types de questions.

\begin{propriete}[les quatre lectures]
Soit $(\Cf)$ la courbe de $f$ et $(\mathscr{C}_{g})$ celle de $g$.
\begin{itemize}[leftmargin=1.4em, itemsep=0.22em]
  \item $f(x) = g(x)$ : abscisses des \emph{points d'intersection} des deux
        courbes.
  \item $f(x) \leq g(x)$ : abscisses des points où $(\Cf)$ est \emph{au-dessous}
        de $(\mathscr{C}_{g})$.
  \item $f(x) = m$ : abscisses des points d'intersection de $(\Cf)$ avec la
        droite \emph{horizontale} d'ordonnée $m$.
  \item $f(x) \leq m$ : abscisses des points où $(\Cf)$ est au-dessous de cette
        horizontale.
\end{itemize}
\end{propriete}

\begin{visualisation}[compter les solutions sans calculer]
\begin{center}
\begin{tikzpicture}
\begin{axis}[
  width = 11cm, height = 5.8cm,
  axis lines = middle, xlabel = {$x$},
  xmin = -3.2, xmax = 3.2, ymin = -3.2, ymax = 5.2,
  xtick = \empty, ytick = {-1, 1, 3}, yticklabels = {$m_3$, $m_2$, $m_1$},
  axis line style = {bmtGris},
  tick label style = {font=\footnotesize, color=bmtGris}, clip = true,
]
  \addplot[bmtIndigo, line width=1.3pt, domain=-2.6:2.6, samples=200] {x^3 - 3*x};
  \addplot[bmtLaterite, dashed, line width=0.85pt, domain=-3:3, samples=2] {3};
  \addplot[bmtVetiver, dashed, line width=0.85pt, domain=-3:3, samples=2] {1};
  \addplot[bmtRaphia, dashed, line width=0.85pt, domain=-3:3, samples=2] {-1};
\end{axis}
\end{tikzpicture}
\end{center}

Trois horizontales, trois réponses différentes pour l'équation $f(x) = m$.
La plus haute ne coupe la courbe qu'une fois : une solution. Celle du milieu
la coupe trois fois : trois solutions. La plus basse aussi. Faites glisser
mentalement l'horizontale de haut en bas et comptez les intersections : vous
lisez d'un coup le nombre de solutions selon $m$, sans résoudre la moindre
équation du troisième degré.
\end{visualisation}

\begin{entrainement}[lecture graphique]
\exo[\niveau{1}] La courbe de $f$ passe par $(0;2)$, admet un maximum $3$ en
$x=1$ et tend vers $-\infty$ aux deux infinis. Combien l'équation
$f(x) = 3$ a-t-elle de solutions ? Et $f(x) = 4$ ?

\exo[\niveau{2}] À partir du tableau de variation de la fonction de l'exemple
mené plus haut, discuter selon $m$ le nombre de solutions de $f(x) = m$.

\exo[\niveau{2}] Deux courbes se coupent en trois points d'abscisses $-2$, $0$
et $3$, et $(\Cf)$ est au-dessus de $(\mathscr{C}_{g})$ sur
$\intervalleo{-2}{0}$ et sur $\left]3 ; +\infty\right[$. Résoudre
graphiquement $f(x) \leq g(x)$.
\end{entrainement}

\begin{corrige}
\textbf{1.} $f(x)=3$ : une solution (le maximum est atteint une fois).
$f(x)=4$ : aucune solution, puisque $4$ dépasse le maximum.

\textbf{2.} Avec $f(0) = -2$ maximum local et $f(2) = 2$ minimum local :
pour $m < -2$ ou $m > 2$, une solution ; pour $m = -2$ ou $m = 2$, deux
solutions ; pour $-2 < m < 2$, aucune solution sur un intervalle et deux sur
l'autre — au total deux.

\textbf{3.} L'ensemble des solutions est
$\left]-\infty ; -2\right] \cup \intervalle{0}{3}$.
\end{corrige}

% ===========================================================================
% ===========================================================================
\begin{annales}
Les cinq études qui suivent sont reprises de sujets réellement posés ; leur
provenance est indiquée à chaque fois. Toutes portent sur des fonctions
polynômes ou rationnelles, et se mènent entièrement avec les outils des trois
premiers chapitres.

\exoann{Baccalauréat probatoire, Cameroun, série D, session 2022 — exercice 4}
Soit $f$ définie sur $\intervallefo{0}{+\infty}$ par
$f(x) = \dfrac{3x}{3+4x}$.
\begin{enumerate}[label=\textbf{\alph*)}, leftmargin=*, itemsep=0.15em]
  \item Dresser le tableau de variation de $f$.
  \item Montrer que la droite d'équation $y = \frac34$ est asymptote à la
        courbe de $f$ et préciser la position de la courbe par rapport à
        cette droite.
  \item Construire la courbe de $f$ dans un repère orthonormé.
\end{enumerate}

\exoann{Baccalauréat probatoire, Cameroun, série IH, session 2023 — problème}
Soit $f$ la fonction définie sur $\intervalle{-4}{1}$ par
$f(x) = 2x^{2}+5x-3$.
\begin{enumerate}[label=\textbf{\alph*)}, leftmargin=*, itemsep=0.15em]
  \item Dresser le tableau de variation de $f$ sur $\intervalle{-4}{1}$.
  \item Déterminer les points d'intersection de la courbe de $f$ avec les axes
        du repère.
  \item Construire cette courbe.
\end{enumerate}

\exoann{Baccalauréat probatoire, Cameroun, séries D et TI, session 2023 — exercice 4}
Soit $f$ définie par $f(x) = x+\dfrac{1}{x+1}$. Dresser le tableau de variation
de $f$, préciser ses deux asymptotes, puis construire sa courbe.

\exoann{Baccalauréat blanc, lycée Philibert Tsiranana, Terminale S, 2022-2023}
Soit $f$ définie sur $\R\setminus\{2\}$ par
$f(x) = \dfrac{x^{2}-3x+4}{x-2}$.
\begin{enumerate}[label=\textbf{\alph*)}, leftmargin=*, itemsep=0.15em]
  \item Déterminer trois réels $a$, $b$ et $c$ tels que
        $f(x) = ax+b+\dfrac{c}{x-2}$.
  \item En déduire que la droite d'équation $y = x-1$ est asymptote oblique à
        la courbe de $f$, et étudier la position relative de la courbe et de
        cette droite.
  \item Dresser le tableau de variation de $f$, puis construire sa courbe et
        ses asymptotes.
\end{enumerate}

\exoann{D'après le concours d'entrée, ENI Fianarantsoa, session 2002-2003 — exercice 1}
Soit $f_{0}$ la fonction définie sur $\R$ par
$f_{0}(x) = \dfrac{x^{2}+3x}{x^{2}+x+3}$ et $(C_{0})$ sa courbe.
\begin{enumerate}[label=\textbf{\alph*)}, leftmargin=*, itemsep=0.15em]
  \item Montrer que $f'_{0}(x)$ est du signe de $-2x^{2}+6x+9$ et dresser le
        tableau de variation de $f_{0}$.
  \item Montrer que $(C_{0})$ coupe son asymptote horizontale en un unique
        point, dont on donnera l'abscisse.
  \item Construire $(C_{0})$ ainsi que son asymptote horizontale.
\end{enumerate}

\exoann{D'après le baccalauréat probatoire, Cameroun, série C, session 2021 — exercice 3}
Soit $f$ définie sur $\R\setminus\{1\}$ par $f(x) = \dfrac{x^{2}-4}{1-x}$, de
courbe $(C)$.
\begin{enumerate}[label=\textbf{\alph*)}, leftmargin=*, itemsep=0.15em]
  \item Déterminer trois réels $a$, $b$ et $c$ tels que
        $f(x) = ax+b+\dfrac{c}{x-1}$.
  \item En déduire que la droite $(T)$ d'équation $y = -x-1$ est asymptote
        oblique à $(C)$, et étudier la position de $(C)$ par rapport à $(T)$.
  \item Construire $(C)$, son asymptote verticale et $(T)$ dans un même
        repère.
\end{enumerate}

\exoann{D'après le baccalauréat probatoire, Cameroun, série C, session 2023 — exercice 4}
Soit $f$ définie par $f(x) = \dfrac{2x^{2}-6x+3}{2x-3}$, de courbe $(C)$.
\begin{enumerate}[label=\textbf{\alph*)}, leftmargin=*, itemsep=0.15em]
  \item Déterminer le réel $c$ tel que
        $f(x) = x-\dfrac32+\dfrac{c}{x-\frac32}$.
  \item Démontrer que la droite $(D)$ d'équation $y = x-\frac32$ est
        asymptote oblique à $(C)$, et préciser les positions relatives de
        $(C)$ et de $(D)$.
  \item Construire $(D)$ et $(C)$.
\end{enumerate}

\exoann{D'après le baccalauréat probatoire, Cameroun, série C, session 2024 — exercice 3}
Soit $f$ définie par $f(x) = 2x+1+\dfrac{1}{2x-3}$, de courbe $(C_{f})$.
\begin{enumerate}[label=\textbf{\alph*)}, leftmargin=*, itemsep=0.15em]
  \item Démontrer que la droite $(D)$ d'équation $y = 2x+1$ est asymptote à
        $(C_{f})$.
  \item Tracer $(C_{f})$ et $(D)$.
  \item Discuter, suivant les valeurs du réel $m$, le nombre de solutions de
        l'équation $f(x) = m$.
\end{enumerate}

\exoann{D'après le baccalauréat probatoire, Cameroun, série C, session 2025 — exercice 2}
Soit $f$ définie par $f(x) = \dfrac{x^{2}-2x+2}{x-1}$.
\begin{enumerate}[label=\textbf{\alph*)}, leftmargin=*, itemsep=0.15em]
  \item Montrer que $f(x) = x-1+\dfrac{1}{x-1}$, puis donner les deux
        asymptotes de la courbe de $f$.
  \item Vérifier que $f'(x) = \dfrac{x^{2}-2x}{(x-1)^{2}}$ et dresser le
        tableau de variation de $f$ sur $\intervalleo{1}{+\infty}$.
  \item Tracer la courbe de $f$ sur $\intervalleo{1}{+\infty}$ avec ses
        asymptotes.
\end{enumerate}

\exoann{D'après le baccalauréat probatoire, Cameroun, série C, session 2022 — exercice 2}
Soit $f$ définie sur $\R$ par
$f(x) = \dfrac{1}{3}x^{3}+\dfrac{3}{2}x^{2}-4x-\dfrac{31}{6}$. Dresser le
tableau de variation de $f$, puis construire sa courbe en plaçant les deux
extremums locaux et le point $(-1\,;0)$.

\medskip
\bmtsource{Les autres énoncés d'annales portant sur l'étude et la
représentation graphique d'une fonction font intervenir le logarithme ou
l'exponentielle. Ils figurent donc aux chapitres 5 et 6, où ils sont signalés
par la mention « étude de fonction ».}
\end{annales}

\begin{corrige}[annales]
\textbf{1. a)} $f'(x) = \dfrac{9}{(3+4x)^{2}} > 0$ : $f$ croît strictement de
$f(0) = 0$ vers $\frac34$ sur $\intervallefo{0}{+\infty}$.
\textbf{b)} $f(x)-\dfrac34 = \dfrac{12x-3(3+4x)}{4(3+4x)}
= \dfrac{-9}{4(3+4x)} < 0$ pour $x \geq 0$ : la courbe reste sous son
asymptote, et l'écart tend vers $0$.
\textbf{c)} La courbe part de l'origine, monte sans jamais atteindre le
palier $y = \frac34$ ; la tangente en $0$ a pour coefficient directeur
$f'(0) = 1$.

\textbf{2. a)} $f'(x) = 4x+5$ s'annule en $-\frac54$. Sur
$\intervalle{-4}{-\frac54}$ la fonction décroît de $f(-4) = 9$ à
$f\left(-\frac54\right) = -\frac{49}{8}$ ; sur
$\intervalle{-\frac54}{1}$ elle croît jusqu'à $f(1) = 4$.
\textbf{b)} $f(0) = -3$ donne le point $(0\,;-3)$ sur l'axe des ordonnées.
L'équation $2x^{2}+5x-3 = 0$ a pour discriminant $49$ et pour racines
$\frac12$ et $-3$ : la courbe coupe l'axe des abscisses en
$\left(\frac12\,;0\right)$ et $(-3\,;0)$.
\textbf{c)} C'est un arc de parabole de sommet
$S\left(-\frac54\,;-\frac{49}{8}\right)$, tracé entre $x = -4$ et $x = 1$.

\textbf{3.} On a vu que
$f'(x) = \dfrac{x(x+2)}{(x+1)^{2}}$. La fonction croît sur
$\intervalleof{-\infty}{-2}$ jusqu'au maximum local $f(-2) = -3$, décroît sur
$\intervallefo{-2}{-1}$ puis sur $\intervalleof{-1}{0}$ jusqu'au minimum local
$f(0) = 1$, et croît ensuite. La droite $x = -1$ est asymptote verticale, la
droite $y = x$ asymptote oblique aux deux infinis ; comme
$f(x)-x = \dfrac{1}{x+1}$, la courbe est au-dessus de cette droite pour
$x > -1$ et en dessous pour $x < -1$. Les deux branches se dessinent de part
et d'autre de l'asymptote verticale.

\textbf{4. a)} La division euclidienne de $x^{2}-3x+4$ par $x-2$ donne
$x^{2}-3x+4 = (x-2)(x-1)+2$, d'où
$f(x) = x-1+\dfrac{2}{x-2}$ : $a = 1$, $b = -1$, $c = 2$.
\textbf{b)} $f(x)-(x-1) = \dfrac{2}{x-2} \to 0$ aux deux infinis : la droite
$y = x-1$ est asymptote oblique. Cette différence est positive pour $x>2$ et
négative pour $x<2$ : la courbe est au-dessus de l'asymptote à droite de $2$,
en dessous à gauche.
\textbf{c)} $f'(x) = 1-\dfrac{2}{(x-2)^{2}} = \dfrac{(x-2)^{2}-2}{(x-2)^{2}}$
s'annule en $2\pm\sqrt2$. La fonction croît jusqu'à
$f\left(2-\sqrt2\right) = 1-2\sqrt2 \approx -1{,}83$, décroît sur
$\intervalle{2-\sqrt2}{2}$ puis sur $\intervalle{2}{2+\sqrt2}$ jusqu'à
$f\left(2+\sqrt2\right) = 1+2\sqrt2 \approx 3{,}83$, et croît ensuite. La
courbe possède deux branches séparées par l'asymptote verticale $x = 2$.

\textbf{5. a)} En dérivant le quotient,
\[
  f'_{0}(x) = \frac{(2x+3)\left(x^{2}+x+3\right)-\left(x^{2}+3x\right)(2x+1)}
  {\left(x^{2}+x+3\right)^{2}}
  = \frac{-2x^{2}+6x+9}{\left(x^{2}+x+3\right)^{2}} .
\]
Le dénominateur étant strictement positif, le signe de $f'_{0}$ est celui du
trinôme $-2x^{2}+6x+9$, de discriminant $108$ et de racines
$\dfrac{3-3\sqrt3}{2} \approx -1{,}10$ et $\dfrac{3+3\sqrt3}{2} \approx 4{,}10$.
La fonction décroît, puis croît entre ces deux valeurs, puis décroît de
nouveau ; les limites en $\pm\infty$ valent $1$, et la droite $y = 1$ est
asymptote horizontale.
\textbf{b)} L'égalité $f_{0}(x) = 1$ équivaut à
$x^{2}+3x = x^{2}+x+3$, c'est-à-dire à $2x = 3$ : la courbe traverse son
asymptote au point d'abscisse $\frac32$, et en ce seul point.
\textbf{c)} La courbe monte depuis $y = 1$ à gauche, plonge jusqu'au minimum
$f_{0}\!\left(\frac{3-3\sqrt3}{2}\right) \approx -0{,}67$, remonte en
franchissant l'asymptote en $x = \frac32$, atteint son maximum
$f_{0}\!\left(\frac{3+3\sqrt3}{2}\right) \approx 1{,}22$, puis redescend
lentement vers $1$.

\textbf{6. a)} On réduit au même dénominateur :
$-x-1+\dfrac{3}{x-1} = \dfrac{-(x+1)(x-1)+3}{x-1}
= \dfrac{-x^{2}+4}{x-1} = \dfrac{x^{2}-4}{1-x}$. Donc $a = -1$, $b = -1$ et
$c = 3$.
\textbf{b)} $f(x)-(-x-1) = \dfrac{3}{x-1}$, qui tend vers $0$ aux deux
infinis : la droite $(T)$ est asymptote oblique. Cette différence est
strictement positive pour $x>1$ et strictement négative pour $x<1$ : la courbe
est au-dessus de $(T)$ à droite de $1$, au-dessous à gauche.
\textbf{c)} La courbe possède deux branches, séparées par l'asymptote
verticale $x = 1$ ; toutes deux décroissent, l'une venant de $+\infty$ et
descendant vers $-\infty$ le long de $x = 1$, l'autre repartant de $+\infty$
au-dessus de $(T)$ pour descendre vers $-\infty$.

\textbf{7. a)} On réduit au même dénominateur :
$x-\dfrac32-\dfrac{3/4}{x-\frac32}
= \dfrac{\left(x-\frac32\right)^{2}-\frac34}{x-\frac32}$. En multipliant
numérateur et dénominateur par $2$, on retrouve
$\dfrac{2x^{2}-6x+3}{2x-3}$. Donc $c = -\dfrac34$.
\textbf{b)} $f(x)-\left(x-\frac32\right) = -\dfrac{3/4}{x-\frac32}$ tend vers
$0$ aux deux infinis : $(D)$ est asymptote oblique. Cette différence est
négative pour $x>\frac32$ et positive pour $x<\frac32$ : la courbe est
au-dessous de $(D)$ à droite de $\frac32$, au-dessus à gauche.
\textbf{c)} $f$ étant strictement croissante sur chacun des deux intervalles,
les deux branches montent, l'une vers l'asymptote verticale par la gauche,
l'autre en repartant de $-\infty$.

\textbf{8. a)} $f(x)-(2x+1) = \dfrac{1}{2x-3}$, qui tend vers $0$ aux deux
infinis : $(D)$ est asymptote à $(C_{f})$.
\textbf{b)} La courbe atteint un maximum local $f(1) = 2$ et un minimum local
$f(2) = 6$, séparés par l'asymptote verticale $x = \frac32$.
\textbf{c)} Sur $\left]-\infty\,;\frac32\right[$, $f$ croît de $-\infty$ à
$f(1) = 2$ puis décroît vers $-\infty$ : elle y prend deux fois toute valeur
strictement inférieure à $2$, une seule fois la valeur $2$, jamais de valeur
supérieure. Sur $\left]\frac32\,;+\infty\right[$, elle décroît de $+\infty$ à
$f(2) = 6$ puis croît vers $+\infty$ : elle y prend deux fois toute valeur
strictement supérieure à $6$, une seule fois la valeur $6$, jamais de valeur
inférieure. Au total :
\[\begin{aligned}
m<2 &: \text{deux solutions} ; \\
  m = 2 &: \text{une} ; \\
  2<m<6 &: \text{aucune} ; \\
  m = 6 &: \text{une} ; \\
  m>6 &: \text{deux.}
\end{aligned}\]

\textbf{9. a)} $x^{2}-2x+2 = (x-1)^{2}+1$, donc
$f(x) = \dfrac{(x-1)^{2}+1}{x-1} = x-1+\dfrac{1}{x-1}$. La droite $x = 1$ est
asymptote verticale, et $f(x)-(x-1) = \frac{1}{x-1}\to0$ montre que
$y = x-1$ est asymptote oblique.
\textbf{b)} $f'(x) = 1-\dfrac{1}{(x-1)^{2}} = \dfrac{(x-1)^{2}-1}{(x-1)^{2}}
= \dfrac{x^{2}-2x}{(x-1)^{2}}$, du signe de $x(x-2)$. Sur
$\intervalleo{1}{+\infty}$ ce produit est négatif avant $2$ et positif après :
$f$ décroît de $+\infty$ à $f(2) = 2$, puis croît vers $+\infty$.
\textbf{c)} La courbe est une branche unique, en forme de creux, tangente
à aucune des deux asymptotes mais s'en approchant de part et d'autre du
minimum $(2\,;2)$.

\textbf{10.} Le tableau reprend le calcul de l'exercice précédent :
croissance sur $\left]-\infty\,;-4\right]$ jusqu'à $f(-4) = \frac{27}{2}$,
décroissance sur $\intervalle{-4}{1}$ jusqu'à $f(1) = -\frac{22}{3}$,
croissance ensuite. La courbe coupe l'axe des abscisses trois fois, dont une
au point $(-1\,;0)$, et présente un point d'inflexion en
$x = -\frac32$, où $f''(x) = 2x+3$ s'annule en changeant de signe.
\end{corrige}

\begin{jecherche}[la boîte de conserve]
Une conserverie de Mahajanga produit des boîtes cylindriques de volume
$500$ cm$^{3}$. On note $r$ le rayon de la base, en centimètres, et $h$ la
hauteur.

\begin{enumerate}[label=\textbf{\arabic*.}, leftmargin=*, itemsep=0.3em]
  \item Exprimer $h$ en fonction de $r$.
  \item Montrer que l'aire totale de métal utilisée, en cm$^{2}$, est
        \[
          A(r) = 2\pi r^{2} + \frac{1000}{r}, \qquad r > 0 .
        \]
  \item Étudier les limites de $A$ en $0^{+}$ et en $+\infty$, puis interpréter
        chacune concrètement.
  \item Déterminer les variations de $A$ et en déduire le rayon qui minimise
        la quantité de métal.
  \item La courbe de $A$ admet-elle une asymptote oblique ? Justifier.
  \item L'entreprise fabrique actuellement des boîtes de rayon $3$ cm. De
        quel pourcentage pourrait-elle réduire sa consommation de métal ?
\end{enumerate}
\end{jecherche}

\begin{corrige}[la boîte de conserve]
\textbf{1.} Le volume vaut $\pi r^{2}h = 500$, donc
$h = \dfrac{500}{\pi r^{2}}$.

\textbf{2.} L'aire totale comprend deux disques et la surface latérale :
\[
  A = 2\pi r^{2} + 2\pi r h = 2\pi r^{2} + 2\pi r \times \frac{500}{\pi r^{2}}
    = 2\pi r^{2} + \frac{1000}{r}.
\]

\textbf{3.} En $0^{+}$, $\dfrac{1000}{r} \to +\infty$ : une boîte très étroite
est très haute, et sa paroi latérale coûte une quantité de métal démesurée. En
$+\infty$, $2\pi r^{2} \to +\infty$ : une boîte très large est très plate, et
ce sont les deux disques qui coûtent cher.

\textbf{4.} $A'(r) = 4\pi r - \dfrac{1000}{r^{2}} = \dfrac{4\pi r^{3}-1000}{r^{2}}$,
du signe de $4\pi r^{3}-1000$. Elle s'annule pour
$r = \sqrt[3]{\dfrac{250}{\pi}} \approx 4{,}30$ cm, en passant du négatif au
positif : c'est un minimum.

\textbf{5.} $\dfrac{A(r)}{r} = 2\pi r + \dfrac{1000}{r^{2}} \to +\infty$ :
aucune asymptote oblique, la courbe admet une branche parabolique de direction
$(Oy)$.

\textbf{6.} $A(3) \approx 389{,}9$ cm$^{2}$ et $A(4{,}30) \approx 348{,}8$
cm$^{2}$. L'économie possible est d'environ
$\dfrac{389{,}9-348{,}8}{389{,}9} \approx 10{,}5\,\%$ du métal.
\end{corrige}


% ===========================================================================
\begin{jemeteste}
Répondre par vrai ou faux, en justifiant en une ligne.

\begin{enumerate}[label=\textbf{\arabic*.}, leftmargin=*, itemsep=0.28em]
  \item Une courbe ne peut pas couper son asymptote oblique.
  \item Si $\limpinf \dfrac{f(x)}{x} = 2$, la droite $y = 2x$ est asymptote à
        la courbe.
  \item Une fonction peut admettre deux asymptotes horizontales.
  \item En un point angulé, la fonction est continue.
  \item Si $f$ est définie sur $\R$ et $\limpinf f(x) = +\infty$, la courbe
        admet une branche parabolique de direction $(Oy)$.
  \item Le tracé d'une courbe doit commencer par le placement des asymptotes.
\end{enumerate}
\end{jemeteste}

\begin{corrige}[je me teste]
\textbf{1.} Faux — elle peut la couper, et même plusieurs fois ; seule la
limite de l'écart compte.
\textbf{2.} Faux — il faut encore que $f(x)-2x$ ait une limite finie.
\textbf{3.} Vrai — une en $+\infty$, une autre en $-\infty$.
\textbf{4.} Vrai — c'est dans la définition.
\textbf{5.} Faux — cela dépend de $\limpinf \frac{f(x)}{x}$ ; une asymptote
oblique est possible.
\textbf{6.} Vrai — c'est ce qui rend le reste du tracé fiable.
\end{corrige}

% ===========================================================================
\begin{commeaubac}
\bmtsource{Exercice au format de l'épreuve}
Soit $f$ la fonction définie sur $\R \setminus \{-1\}$ par
$f(x) = \dfrac{2x^{2}+3x-1}{x+1}$, de courbe $(\Cf)$.

\begin{enumerate}[label=\textbf{\arabic*.}, leftmargin=*, itemsep=0.25em]
  \item Déterminer trois réels $a$, $b$ et $c$ tels que
        $f(x) = ax+b+\dfrac{c}{x+1}$. \hfill\textup{(0,75 pt)}
  \item En déduire les asymptotes de $(\Cf)$ et la position de la courbe par
        rapport à l'asymptote oblique. \hfill\textup{(1,25 pt)}
  \item Calculer $f'(x)$, étudier son signe et dresser le tableau de variation
        de $f$. \hfill\textup{(1,5 pt)}
  \item Montrer que le point $I(-1;-1)$ est centre de symétrie de $(\Cf)$.
        \hfill\textup{(1 pt)}
  \item Construire $(\Cf)$ et ses asymptotes. \hfill\textup{(1,5 pt)}
  \item Discuter graphiquement, selon le réel $m$, le nombre de solutions de
        l'équation $f(x) = m$. \hfill\textup{(1 pt)}
\end{enumerate}
\end{commeaubac}

\begin{corrige}[exercice au format de l'épreuve]
\textbf{1.} La division euclidienne donne $f(x) = 2x+1-\dfrac{2}{x+1}$, donc
$a=2$, $b=1$, $c=-2$.

\textbf{2.} $x = -1$ est asymptote verticale ; $y = 2x+1$ est asymptote
oblique aux deux infinis. Le signe de $f(x)-(2x+1) = \dfrac{-2}{x+1}$ est
l'opposé de celui de $x+1$ : la courbe est au-dessous de l'asymptote pour
$x > -1$, au-dessus pour $x < -1$.

\textbf{3.} $f'(x) = 2 + \dfrac{2}{(x+1)^{2}} > 0$ : $f$ est strictement
croissante sur chacun des deux intervalles, de $-\infty$ à $-\infty$ puis de
$+\infty$ à $+\infty$.

\textbf{4.} On vérifie que $f(-1+h) + f(-1-h) = -2$ pour tout $h \neq 0$ :
en effet $f(-1+h) = 2h-1-\dfrac{2}{h}$ et $f(-1-h) = -2h-1+\dfrac{2}{h}$. La
demi-somme vaut $-1$, ordonnée de $I$ : le point $I$ est centre de symétrie.

\textbf{5.} La courbe possède deux branches, chacune parcourant $\R$ tout
entier de façon strictement croissante. Pour tout réel $m$, l'équation
$f(x)=m$ admet donc exactement deux solutions, une sur chaque branche.
\end{corrige}

% ===========================================================================
\begin{concours}
\bmtsource{D'après le concours d'entrée de l'ENI, Fianarantsoa}
Soit $m$ un paramètre réel et $f_{m}$ la fonction définie par
\[
  f_{m}(x) = \frac{x^{2}+3x+4m}{x^{2}+(5m+1)x+3},
\]
de courbe $(\mathscr{C}_{m})$.

\begin{enumerate}[label=\textbf{\arabic*.}, leftmargin=*, itemsep=0.25em]
  \item Montrer que toutes les courbes $(\mathscr{C}_{m})$ passent par trois
        points fixes, dont on déterminera les coordonnées.
  \item Déterminer $m$ pour que le point d'intersection de $(\mathscr{C}_{m})$
        avec son asymptote parallèle à $(x'x)$ ait pour abscisse
        $\dfrac{3}{2}$.
  \item Construire la courbe $(\mathscr{C}_{0})$.
\end{enumerate}
\end{concours}

\begin{corrige}[vers le concours]
\textbf{1.} Un point $(x_{0};y_{0})$ appartient à toutes les courbes lorsque
$y_{0}\bigl(x_{0}^{2}+(5m+1)x_{0}+3\bigr) = x_{0}^{2}+3x_{0}+4m$ pour
\emph{tout} $m$. En regroupant les termes en $m$ :
\[
  m\bigl(5x_{0}y_{0}-4\bigr)
  + \bigl(y_{0}x_{0}^{2}+y_{0}x_{0}+3y_{0}-x_{0}^{2}-3x_{0}\bigr) = 0 .
\]
Les deux parenthèses doivent s'annuler séparément. La première donne
$y_{0} = \dfrac{4}{5x_{0}}$ ; en la reportant dans la seconde on obtient
$5x_{0}^{3}+11x_{0}^{2}-4x_{0}-12 = 0$, dont $1$ est racine évidente, d'où la
factorisation $(x_{0}-1)(5x_{0}^{2}+16x_{0}+12) = 0$ et les racines
$x_{0} \in \left\{1\,;-\tfrac{6}{5}\,;-2\right\}$. Les trois points fixes sont
donc
\[
  \left(1\,;\tfrac45\right), \qquad
  \left(-\tfrac65\,;-\tfrac23\right), \qquad
  \left(-2\,;-\tfrac25\right).
\]

\textbf{2.} Numérateur et dénominateur ayant même degré et même coefficient
dominant, l'asymptote horizontale a pour équation $y = 1$. L'équation
$f_{m}(x) = 1$ donne $3x+4m = (5m+1)x+3$, soit $x = \dfrac{3-4m}{2-5m}$.
Imposer $x = \dfrac32$ conduit à $2(3-4m) = 3(2-5m)$, d'où $7m = 0$ et
$m = 0$.

\textbf{3.} Pour $m = 0$, $f_{0}(x) = \dfrac{x^{2}+3x}{x^{2}+x+3}$. Le
dénominateur ne s'annule jamais : aucune asymptote verticale, seule
l'asymptote horizontale $y=1$ subsiste. La courbe coupe l'axe des abscisses en
$0$ et en $-3$.
\end{corrige}


% ===========================================================================
\begin{devoir}[2 heures]
Trois exercices indépendants, notés sur $20$. Le devoir porte sur ce chapitre
et sur les deux qui le précèdent. Les énoncés sont repris de sessions réelles
et assemblés ici en épreuve ; leur provenance est indiquée à chaque fois.

\exods{1}{D'après le baccalauréat probatoire, Cameroun, série C, session 2022 — exercice 2}{5 points}
Soit $f$ la fonction définie sur $\R$ par
$f(x) = \dfrac{1}{3}x^{3}+\dfrac{3}{2}x^{2}-4x-\dfrac{31}{6}$.
\begin{enumerate}[label=\textbf{\arabic*.}, leftmargin=*, itemsep=0.2em]
  \item Vérifier que la courbe de $f$ passe par le point $(-1\,;0)$.
        \hfill\textup{(0,5 pt)}
  \item Calculer $f'(x)$ et étudier son signe. \hfill\textup{(1,5 pt)}
  \item Dresser le tableau de variation de $f$. \hfill\textup{(1,5 pt)}
  \item Démontrer que l'équation $f(x) = 0$ admet exactement trois solutions
        réelles. \hfill\textup{(1,5 pt)}
\end{enumerate}

\exods{2}{D'après le baccalauréat probatoire, Cameroun, série C, session 2020 — exercice 3, partie II}{7 points}
Soit $f$ la fonction définie sur $\R\setminus\{-1\}$ par
$f(x) = x+2+\dfrac{4}{x+1}$, de courbe $(C)$.
\begin{enumerate}[label=\textbf{\arabic*.}, leftmargin=*, itemsep=0.2em]
  \item Calculer $f'(x)$, étudier son signe et dresser le tableau de
        variation de $f$. \hfill\textup{(2,5 pts)}
  \item Montrer que la droite $(D)$ d'équation $y = x+2$ est asymptote à
        $(C)$, et préciser la position de $(C)$ par rapport à $(D)$.
        \hfill\textup{(2 pts)}
  \item Déterminer les points d'intersection de $(C)$ avec les axes du
        repère. \hfill\textup{(1 pt)}
  \item Construire $(C)$, $(D)$ et l'asymptote verticale dans un même repère.
        \hfill\textup{(1,5 pt)}
\end{enumerate}

\exods{3}{D'après le baccalauréat probatoire, Cameroun, série C, session 2024 — exercice 3}{8 points}
Soit $f$ la fonction définie par $f(x) = 2x+1+\dfrac{1}{2x-3}$, de courbe
$(C_{f})$.
\begin{enumerate}[label=\textbf{\arabic*.}, leftmargin=*, itemsep=0.2em]
  \item Donner l'ensemble de définition de $f$ et ses limites aux bornes de
        cet ensemble. \hfill\textup{(1,5 pt)}
  \item Calculer $f'(x)$, résoudre l'inéquation $f'(x) \leq 0$, puis dresser
        le tableau de variation de $f$. \hfill\textup{(2 pts)}
  \item Montrer que la droite $(D)$ d'équation $y = 2x+1$ est asymptote à
        $(C_{f})$ et préciser les positions relatives. \hfill\textup{(1,5 pt)}
  \item Démontrer que le point $\Omega\left(\frac32\,;4\right)$ est centre de
        symétrie de $(C_{f})$. \hfill\textup{(1,5 pt)}
  \item Construire $(C_{f})$, puis discuter suivant les valeurs du réel $m$
        le nombre de solutions de l'équation $f(x) = m$.
        \hfill\textup{(1,5 pt)}
\end{enumerate}
\end{devoir}

\begin{corrige}[devoir surveillé]
\textbf{Exercice 1. 1.} $f(-1) = -\frac13+\frac32+4-\frac{31}{6}
= \frac{-2+9+24-31}{6} = 0$.

\textbf{2.} $f'(x) = x^{2}+3x-4 = (x+4)(x-1)$ : positif sur
$\left]-\infty\,;-4\right]$, négatif sur $\intervalle{-4}{1}$, positif sur
$\intervallefo{1}{+\infty}$.

\textbf{3.} $f$ croît de $-\infty$ jusqu'à $f(-4) = \frac{27}{2}$, décroît
jusqu'à $f(1) = -\frac{22}{3}$, puis croît vers $+\infty$.

\textbf{4.} Sur chacun des trois intervalles de monotonie, $f$ est continue
et strictement monotone, et ses valeurs aux extrémités sont de signes
contraires : $-\infty$ puis $\frac{27}{2}>0$ ; $\frac{27}{2}>0$ puis
$-\frac{22}{3}<0$ ; $-\frac{22}{3}<0$ puis $+\infty$. Le théorème des valeurs
intermédiaires fournit donc une racine et une seule dans chacun, soit trois
racines en tout. La question \textbf{1} en donne une : $-1$, qui appartient à
$\intervalle{-4}{1}$.

\medskip
\textbf{Exercice 2. 1.} $f'(x) = 1-\dfrac{4}{(x+1)^{2}}
= \dfrac{(x+1)^{2}-4}{(x+1)^{2}} = \dfrac{(x+3)(x-1)}{(x+1)^{2}}$, du signe de
$(x+3)(x-1)$. La fonction croît sur $\left]-\infty\,;-3\right]$ jusqu'au
maximum local $f(-3) = -3+2-2 = -3$, décroît sur $\intervallefo{-3}{-1}$ puis
sur $\intervalleof{-1}{1}$ jusqu'au minimum local $f(1) = 1+2+2 = 5$, et croît
ensuite.

\textbf{2.} $f(x)-(x+2) = \dfrac{4}{x+1}$, qui tend vers $0$ aux deux
infinis : $(D)$ est asymptote oblique. Cette différence est strictement
positive pour $x>-1$ et strictement négative pour $x<-1$ : la courbe est
au-dessus de $(D)$ à droite de l'asymptote verticale, au-dessous à gauche.

\textbf{3.} $f(0) = 6$ : la courbe coupe l'axe des ordonnées en $(0\,;6)$.
Pour l'axe des abscisses, $f(x) = 0$ équivaut, après multiplication par
$x+1 \neq 0$, à $x^{2}+3x+6 = 0$, dont le discriminant vaut $9-24 = -15<0$ :
il n'y a aucun point d'intersection avec l'axe des abscisses.

\textbf{4.} La branche de gauche monte jusqu'au point $(-3\,;-3)$ — où la
tangente est horizontale — puis plonge le long de l'asymptote $x = -1$ ; celle
de droite descend depuis $+\infty$ jusqu'au point $(1\,;5)$, puis remonte en
longeant $(D)$ par-dessus.

\medskip
\textbf{Exercice 3. 1.} $\Df = \R\setminus\left\{\frac32\right\}$. Aux
infinis, $\frac{1}{2x-3}\to0$ et $f(x)$ se comporte comme $2x+1$ : les limites
sont $-\infty$ et $+\infty$. En $\frac32^{-}$, $\frac{1}{2x-3}\to-\infty$ donc
$f(x)\to-\infty$ ; en $\frac32^{+}$, $f(x)\to+\infty$.

\textbf{2.} $f'(x) = 2-\dfrac{2}{(2x-3)^{2}}$, et $f'(x)\leq0$ équivaut à
$(2x-3)^{2}\leq1$, c'est-à-dire à $1 \leq x \leq 2$, la valeur $\frac32$ étant
exclue. La fonction croît sur $\left]-\infty\,;1\right]$ jusqu'au maximum
local $f(1) = 2$, décroît jusqu'à l'asymptote, décroît encore jusqu'au minimum
local $f(2) = 6$, puis croît.

\textbf{3.} $f(x)-(2x+1) = \dfrac{1}{2x-3}\to0$ aux deux infinis : $(D)$ est
asymptote. La différence est positive à droite de $\frac32$, négative à
gauche.

\textbf{4.} Pour $h \neq 0$ tel que $\frac32\pm h$ soit défini,
\[\begin{aligned}
&f\!\left(\tfrac32+h\right)+f\!\left(\tfrac32-h\right)\\
  &\quad = \left(4+2h+\frac{1}{2h}\right)+\left(4-2h-\frac{1}{2h}\right) = 8,
\end{aligned}\]
soit deux fois l'ordonnée de $\Omega$ : le point $\Omega$ est centre de
symétrie.

\textbf{5.} Sur $\left]-\infty\,;\frac32\right[$ la fonction monte de
$-\infty$ à $2$ puis redescend vers $-\infty$ ; sur
$\left]\frac32\,;+\infty\right[$ elle descend de $+\infty$ à $6$ puis remonte.
D'où la discussion :
\[\begin{aligned}
m<2 &: \text{deux solutions} ; \\
  m=2 &: \text{une} ; \\
  2<m<6 &: \text{aucune} ; \\
  m=6 &: \text{une} ; \\
  m>6 &: \text{deux.}
\end{aligned}\]
\end{corrige}
